
En
effet ces conditions sont suffisantes car déjà 4) et 5) entraînent l'effectivité. Réciproquement soit @@M0@@ inversible sur @@M1@@ tel qu'il existe un isomorphisme @@M2@@ compatible avec la donnée de recollement. On peut recouvrir @@M3@@ par des ouverts @@M4@@, au-dessus desquels @@M5@@ est isomorphe à @@M6@@~; donc la donnée de descente sur @@M7@@ restreinte au-dessus de @@M8@@ est isomorphe à la donnée de descente triviale sur @@M9@@. Enfin les conditions 1)-5) sont de nature locale sur @@M10@@ et elles sont vérifiées d'après les hypothèses de 4.1 pour @@M11@@~; elles sont donc vérifiées pour @@M12@@.

Il reste donc à représenter successivement chacune des conditions 1) à 5) par des monomorphismes de présentation finie. Pour 1) et 2) on applique le

\textbf{Lemme 4.4.} \uline{Soient} @@M13@@ \uline{un schéma,} @@M14@@ \uline{et} @@M15@@ \uline{des morphismes propres de présentation finie,} @@M16@@ \uline{un} @@M17@@\uline{-morphisme. Soit} @@M18@@ \uline{un faisceau sur} @@M19@@ \uline{de présentation finie et} plat \uline{sur} @@M20@@ \uline{et, pour un entier} @@M21@@\uline{, soit} @@M22@@ \uline{le sous-foncteur de} @@M23@@ \uline{qui rend} @@M24@@ \uline{plat} \uline{sur} @@M25@@ \uline{pour tout} @@M26@@ \uline{(précisément,} @@M27@@ \uline{si} @@M28@@ \uline{est} \uline{plat} \uline{pour tout} @@M29@@, @@M30@@ \uline{sinon). Alors} @@M31@@ \uline{est représentable par un monomorphisme surjectif de présentation finie, et alors pour} @@M32@@ \uline{les} @@M33@@ \uline{commutent aux changements de base} @@M34@@\uline{. De plus} @@M35@@ \uline{coïncide avec le foncteur} @@M36@@ \uline{qui rend exacts les foncteurs}
@@M37@@
\uline{où} @@M38@@ \uline{parcourt les} @@M39@@\uline{-modules quasi-cohérents pour tout ouvert affine} @@M40@@ \uline{de} @@M41@@.

En effet l'assertion est locale sur @@M42@@~; on peut donc supposer @@M43@@ quasi-compact. Alors la dimension relative de @@M44@@ est bornée, soit par @@M45@@, et pour @@M46@@, @@M47@@. L'assertion est donc vraie pour @@M48@@~;
